% mooie C++
% http://archive.cs.uu.nl/mirror/CTAN/macros/latex/contrib/relsize/relsize-doc.pdf
\usepackage{relsize}
\iftoggle{charter}{
    \newcommand{\Cpp}{C\nolinebreak\raisebox{.2ex}[0pt][0pt]{\smaller{+\nolinebreak+}}}
}{
\newcommand{\Cpp}{C\nolinebreak\raisebox{.3ex}[0pt][0pt]{\smaller{+\nolinebreak+}}}
}

% define kleuren die in listings gebruikt worden
\definecolor{cgreen}{RGB}{0,128,0}
\definecolor{cdarkred}{RGB}{163,21,21}
\definecolor{cbrown}{RGB}{121,37,0}
\definecolor{cpurple}{RGB}{97,63,153}
\definecolor{cblue}{RGB}{0,0,255}

\definecolor{mgreen}{RGB}{33,138,77}
\definecolor{mpurple}{RGB}{159,31,242}

\definecolor{agreen}{RGB}{0,128,0}
\definecolor{ablue}{RGB}{0,0,255}

% http://archive.cs.uu.nl/mirror/CTAN/macros/latex/contrib/listings/listings.pdf
\usepackage{listings}
\lstset{
    inputencoding={utf8}, 
    literate=
        {ä}{{\"a}}1 
        {ë}{{\"e}}1 
        {ï}{{\"i}}1 
        {ö}{{\"o}}1
        {ü}{{\"u}}1
        {á}{{\'a}}1
        {é}{{\'e}}1
        {í}{{\'i}}1
        {ó}{{\'o}}1
        {ú}{{\'u}}1
        {à}{{\`a}}1
        {è}{{\`e}}1
        {ì}{{\`i}}1
        {ò}{{\`o}}1
        {ù}{{\`u}}1
        {â}{{\^a}}1
        {ê}{{\^e}}1
        {î}{{\^i}}1
        {ô}{{\^o}}1
        {û}{{\^u}}1
        {©}{{\copyright}}1,
    basicstyle=\ttfamily\small,
    tabsize=4,
    showstringspaces=false,
    upquote=true,
    numbers=none,
    aboveskip=.1\baselineskip,
    belowskip=0pt,
    abovecaptionskip=.6\baselineskip,
    belowcaptionskip=0pt,
    includerangemarker=false,
    prebreak={\raisebox{0ex}[0ex][0ex]{\textcolor{black}{\ensuremath{\hookleftarrow}}}}, 
    postbreak={\raisebox{0ex}[0ex][0ex]{\textcolor{black}{\ensuremath{\space\hookrightarrow\space}}}}, 
    breaklines=true, 
    breakatwhitespace=true, 
    breakindent=0.1em,
    numberbychapter=true,
    captionpos=b,
}

\lstdefinestyle{linenumbers}{
    numbers=left, 
    stepnumber=1, 
    numberstyle=\tiny\textcolor{red}, 
    numbersep=1em
}

% C
\lstdefinestyle{cstyle}{
    language=C,
    morekeywords={interrupt},
    stringstyle=\color{cdarkred},
    keywordstyle=\color{cblue},
    identifierstyle=,
    commentstyle=\color{cgreen},
}

% C++
\lstdefinestyle{cppstyle}{
    language=C++,
    morekeywords={alignas,alignof,char16_t,char32_t,concept,constexpr,decltype,noexcept,nullptr,requires,static_assert,thread_local,override,final},
    stringstyle=\color{cdarkred},
    keywordstyle=\color{cblue},
    identifierstyle=,
    commentstyle=\color{cgreen},
}

% MATLAB
\lstdefinestyle{mstyle}{
    language=Matlab,
    morekeywords={},
    stringstyle=\color{mpurple},
    keywordstyle=,
    identifierstyle=,
    commentstyle=\color{mgreen},
}

% AVR assembler
\lstdefinelanguage[AVR]{Assembler}{
    morekeywords=[1]{
        adc,add,adiw,andi,asr,bclr,bld,brbc,brbs,brcc,brcs,break,breq,brge,brhc,brhs,brid,brie,brlo,brlt,brmi,brne,brpl,brsh,brtc,brts,brvc,brvs,bset,bst,cbi,cbr,clc,clh,cli,cln,clr,cls,clt,clv,clz,com,count,cp,cpc,cpi,cpse,dec,eicall,eijmp,elpm,end,eor,espm,fmul,fmuls,fmulsu,in,inc,jmp,ld,ldd,ldi,lds,lsl,lsr,mov,movw,neg,nop,ori,out,pop,push,rcall,ret,reti,rjmp,rol,sbc,sbci,sbi,sbic,sbis,sbiw,sbr,sbrc,sbrs,sec,seh,sei,sen,ser,ses,sev,sez,sleep,st,std,sts,sub,subi,swap,tst,wdr
    },
    morekeywords=[2]{
        .byte,.cseg,.csegsize,.db,.device,.dseg,.dw,.endm,.equ,.eseg,.exit,.include,.includepath,.listmac,.macro,.nolist,.org,.set,.define,.undef,.ifdef,.ifndef,.else,.elseif,.elif,.message,.warning,.error,byte1,byte2,byte3,byte4,exp2,high,hwrd,log2,low,lwrd,page
    },
    alsoletter=.,
    sensitive=false,
    morecomment=[l]{\#},
    morecomment=[l]{;},
    morestring=[b]{"},
    morestring=[b]{'}
}[keywords,comments,strings]

\lstdefinestyle{astyle}{
    language=[AVR]Assembler,
    stringstyle=,
    keywordstyle=[1]\color{ablue},
    keywordstyle=[2],
    identifierstyle=,
    commentstyle=\color{agreen},  
}

% Commando's \floatlstinput en \lstinput lezen code uit (een deel van) een bestand en plaatsen dit in een listing.
% #1 = (optional) extra key=value paren voor lstinputlisting b.v.
%      style = cstyle, cppstyle, mstyle of astyle (C, C++, MATLAB, AVR assembler)
%      style = linenumbers
%      linerange={1-6, 13-27}
% #2 = filenaam (moet in het directory ./progs staan)
% #3 = label=lst:#3
% #4 = titel
\newcommand{\floatlstinput}[4][]{%
\lstinputlisting[
    belowskip=-.2\baselineskip,
    rangebeginprefix=//>,% Alleen nuttig voor C en C++ 
    rangeendprefix=//<,% Alleen nuttig voor C en C++ 
    float=!htbp, 
    label=lst:#3, 
    caption={#4},%    
    #1
]{progs/#2}%
}

% #1 = (optional) extra key=value paren voor lstinputlisting b.v.
%      style = cstyle, cppstyle, mstyle of astyle (C, C++, MATLAB, AVR assembler)
%      style = linenumbers
%      linerange={1-6, 13-27}
% #2 = filenaam (moet in het directory ./progs staan)
% #3 = label=lst:#3
% #4 = titel
\newcommand{\lstinput}[4][]{%
\lstinputlisting[
    aboveskip=.8\baselineskip,
    belowskip=.4\baselineskip,
    rangebeginprefix=//>,% Alleen nuttig voor C en C++ 
    rangeendprefix=//<,% Alleen nuttig voor C en C++ 
    label=lst:#3, 
    caption={#4},%    
    #1
]{progs/#2}%
}

% Commando \lstinputfragment leest code uit (een deel van) een bestand en plaatst dit in een codefragment.
% #1 = (optional) extra key=value paren voor lstinputlisting b.v.
%      style = cstyle, cppstyle, mstyle of astyle (C, C++, MATLAB, AVR assembler)
%      style = linenumbers
%      linerange={1-6, 13-27}
% #2 = filenaam (moet in het directory ./progs staan)
% geen titel, geen label
\newcommand{\lstinputfragment}[2][]{%
\lstinputlisting[
    aboveskip=.8\baselineskip,
    belowskip=.4\baselineskip,
    rangebeginprefix=//>,% Alleen nuttig voor C en C++ 
    rangeendprefix=//<,% Alleen nuttig voor C en C++ 
    #1
]{progs/#2}%
}

% De code in de environments floatlst en lst wordt in een listing geplaatst.
% #1 = (optional) extra key=value paren voor lstset b.v.
%      style = cstyle, cppstyle, mstyle, astyle (C, C++, MATLAB, AVR assembler)
%      style = linenumbers
% #2 = label=lst:#2
% #3 = titel
\lstnewenvironment{floatlst}[3][]{%
    \lstset{
        belowskip=-.2\baselineskip,
        float=!htbp, 
        label=lst:#2, 
        caption={#3},%   
        #1
    }%
}{%
}

% #1 = (optional) extra key=value paren voor lstset b.v.
%      style = cstyle, cppstyle, mstyle, astyle (C, C++, MATLAB, AVR assembler)
%      style = linenumbers
% #2 = label=lst:#2
% #3 = titel
\lstnewenvironment{lst}[3][]{%
    \lstset{
        aboveskip=.8\baselineskip,
        belowskip=.4\baselineskip,
        label=lst:#2, 
        caption={#3},%    
        #1
    }%
}{%
}

% De code in de environments lstfragment wordt in een codefragment geplaatst.
% #1 = (optional) extra key=value paren voor lstset b.v.
%      style = cstyle, cppstyle, mstyle, astyle (C, C++, MATLAB, AVR assembler)
%      style = linenumbers
% geen titel, geen label
\lstnewenvironment{lstfragment}[1][]{%
    \lstset{
        aboveskip=.4\baselineskip,%
        #1
    }%
}{%
}

% \intextcodefont is nodig om \small te verwijderen
\newcommand{\intextcodefont}{\ttfamily}
\lstdefinestyle{lstinlinestyle}{
    basicstyle=\intextcodefont, 
    prebreak={}, 
    postbreak={}, 
    breakindent=0pt,
}

% Commando's \code, \ccode, \cppcode, \mcode en \acode kunnen gebruikt worden om inline code te definiëren.
% #1 = (optional) extra key=value paren voor lstinline.
% #2 = inline code
\newcommand{\code}[2][]{%
    \lstinline[
        style=lstinlinestyle,%
        #1
    ]@#2@%
}

\newcommand{\ccode}[2][]{%
    \lstinline[
        style=cstyle,
        style=lstinlinestyle,%
        #1
    ]@#2@%
}

\newcommand{\cppcode}[2][]{%
    \lstinline[
        style=cppstyle,
        style=lstinlinestyle,%
        #1
    ]@#2@%
}

\newcommand{\mcode}[2][]{%
    \lstinline[
        style=mstyle,
        style=lstinlinestyle,%
        #1
    ]@#2@%
}

\newcommand{\acode}[2][]{%
    \lstinline[
        style=astyle,
        style=lstinlinestyle,%
        #1
    ]@#2@%
}

% met \hcode kun je ercoor zorgen dat lange identifiers afgebroken kunnen worden. Mogelijke afbreekpunten met \\- markeren.
\newcommand{\hcode}[2][]{%
    \lstinline[
        style=lstinlinestyle,
        literate={\\-}{}{0\discretionary{-}{}{}},%
        #1
    ]@#2@%
}

% handige afkorting |bla| voor \ccode{bla} maar pas op: werkt niet in footnotes, captions etc.
% werkt niet als je | in latex code gebruikt (b.v. in math).
% vandaar environment van gemaakt.
\newenvironment{clstShortInline}{%
    \lstMakeShortInline[
        style=cstyle,
        style=lstinlinestyle,
    ]{|}%
}{%
    \lstDeleteShortInline{|}%
}

\newenvironment{cpplstShortInline}{%
    \lstMakeShortInline[
        style=cppstyle,
        style=lstinlinestyle,
    ]{|}%
}{%
    \lstDeleteShortInline{|}%
}

\newenvironment{mlstShortInline}{%
    \lstMakeShortInline[
        style=mstyle,
        style=lstinlinestyle,
    ]{|}%
}{%
    \lstDeleteShortInline{|}%
}

\newenvironment{alstShortInline}{%
    \lstMakeShortInline[
        style=astyle,
        style=lstinlinestyle,
    ]{|}%
}{%
    \lstDeleteShortInline{|}%
}

% gebruik voor opmaak van console output
% http://ftp.snt.utwente.nl/pub/software/tex/macros/latex/contrib/fancyvrb/fancyvrb.pdf
% nodig om underline voor input te kunnen gebruiken: \cinput{...}
\usepackage{fancyvrb}
\definecolor{darkgreen}{RGB}{0,100,0}
\newcommand{\cinput}[1]{\color{darkgreen}{\underline{#1}}}
\newenvironment{coutput}{%
    \Verbatim[commandchars=\\\{\}]%
    }{%
    \endVerbatim%
}

%http://ctan.cs.uu.nl/macros/generic/xstring/xstring_doc_en.pdf
\usepackage{xstring}
\makeatletter
\newrobustcmd{\proglink}[1]{%
    \normalexpandarg%
    \StrDel{#1}{\\-}[\bare@link]%
    \href{\progsurl/\bare@link}{\hcode{#1}}%
}
\makeatother
